\chapter{ICMP, ping e traceroute}

% ══════════════════════════════════════════════════════════════
\section{Il protocollo ICMP}
% ══════════════════════════════════════════════════════════════

\dfn{ICMP}{
    L'\textbf{Internet Control Message Protocol} (ICMP) è un protocollo diagnostico e di controllo usato da host e
    router per scambiarsi informazioni sul livello di rete: segnalazioni di errore (destinazione irraggiungibile,
    TTL scaduto) e richieste di verifica (\emph{echo}). I messaggi ICMP viaggiano come \textbf{payload di datagrammi
    IP} (campo protocollo = 1), ma ICMP è considerato parte del livello di rete.
}

\begin{figure}[H]
\centering
\begin{tikzpicture}[every node/.style={font=\scriptsize\sffamily}]
  \node[field, fill=lnet, minimum width=2.4cm] (ip) at (0,0) {header IP (protocollo = 1)};
  \node[field, fill=netred!25, minimum width=1cm, right=0pt of ip] (ty) {tipo};
  \node[field, fill=netred!25, minimum width=1cm, right=0pt of ty] (co) {codice};
  \node[field, fill=netred!25, minimum width=1.4cm, right=0pt of co] (ck) {checksum};
  \node[field, fill=netred!10, minimum width=5.2cm, right=0pt of ck] {resto del messaggio (dipende dal tipo)};
  \draw[decorate, decoration={brace, amplitude=5pt}] ([yshift=2pt]ty.north west) -- ([yshift=2pt]12.2,0.33) node[midway, above=6pt] {messaggio ICMP, payload del datagramma IP};
\end{tikzpicture}
\caption{Un messaggio ICMP dentro un datagramma IP.}
\end{figure}

\begin{table}[H]
\centering
\begin{tabular}{ccl}
\toprule
\textbf{Tipo} & \textbf{Codice} & \textbf{Significato} \\
\midrule
0 & 0 & echo reply (risposta a ping) \\
3 & 0 & destinazione irraggiungibile: rete irraggiungibile \\
3 & 1 & destinazione irraggiungibile: host irraggiungibile \\
3 & 3 & destinazione irraggiungibile: \textbf{porta} irraggiungibile \\
3 & 4 & frammentazione necessaria ma vietata (``packet too big'') \\
8 & 0 & echo request (richiesta di ping) \\
11 & 0 & \textbf{TTL scaduto} (\emph{time exceeded}) \\
12 & 0 & header IP errato \\
\bottomrule
\end{tabular}
\caption{Alcuni tipi di messaggi ICMP.}
\end{table}

% ══════════════════════════════════════════════════════════════
\section{ping}
% ══════════════════════════════════════════════════════════════

ICMP è un protocollo con cui, sostanzialmente, un client invia pacchetti a un server che risponde: dai
\textbf{timestamp} si ricava il tempo di risposta, e si verifica anche se ci sono \textbf{perdite} di pacchetti.

\dfn{ping}{
    \textbf{ping} (\emph{Packet Internet Groper}) invia a un host messaggi ICMP \textbf{echo request} (tipo 8), a cui l'host
    risponde automaticamente con \textbf{echo reply} (tipo 0). Per ogni coppia mostra il tempo trascorso, cioè
    l'\textbf{RTT}, e alla fine le statistiche sulle perdite.
}

\begin{lstlisting}[language=bash]
$ ping iana.org
PING iana.org (192.0.43.8) 56(84) bytes of data.
64 bytes from 192.0.43.8: icmp_seq=1 ttl=51 time=178 ms
64 bytes from 192.0.43.8: icmp_seq=2 ttl=51 time=176 ms
...
--- iana.org ping statistics ---
4 packets transmitted, 4 received, 0% packet loss
\end{lstlisting}

\warning{ping non usa il livello di trasporto}{
    ping non usa né TCP né UDP: costruisce e riceve direttamente pacchetti ICMP usando un \emph{raw socket}, un socket che
    permette a un programma di scrivere pacchetti a livello di rete. Per questo non ha numeri di porta; al loro posto usa un
    \textbf{identificatore} e un \textbf{numero di sequenza} nel messaggio ICMP, per associare ogni risposta alla sua richiesta.
}

\subsection{Cosa osservare in Wireshark}

\begin{esercizio}{ping con Wireshark}
Si avvia la cattura, si esegue \texttt{ping iana.org} e si filtra con \texttt{icmp}.
\begin{enumerate}
    \item Individuare i pacchetti ICMP \textbf{echo request} inviati dal proprio computer, annotando indirizzi IP sorgente e destinazione.
    \item Trovare i corrispondenti \textbf{echo reply} dal server.
    \item Esaminare un echo request:
    \begin{itemize}
      \item nella sezione IP, osservare il campo \textbf{TTL};
      \item nella sezione ICMP, trovare il campo \textbf{Type} (8 per echo request);
      \item annotare \textbf{identificatore} e \textbf{numero di sequenza};
      \item misurare l'\textbf{RTT} tra echo request ed echo reply.
    \end{itemize}
\end{enumerate}
\end{esercizio}

% ══════════════════════════════════════════════════════════════
\section{traceroute}
% ══════════════════════════════════════════════════════════════

\textbf{traceroute} scopre la sequenza di router attraversati per raggiungere una destinazione. Sfrutta in modo
astuto il campo \textbf{TTL} del datagramma IP.

\subsection{Il meccanismo}

\begin{itemize}
    \item Ogni volta che un pacchetto attraversa un router (un \emph{hop}), il suo TTL diminuisce di uno.
    \item traceroute usa deliberatamente un TTL basso, \textbf{1} per cominciare, così che il pacchetto ``scada'' al
          \textbf{primo router}.
    \item Il router che riceve il pacchetto con TTL scaduto lo scarta e invia al mittente un messaggio ICMP
          \textbf{Time Exceeded} (tipo 11). L'indirizzo IP sorgente di quel messaggio è l'indirizzo del router!
    \item traceroute invia quindi una sequenza di pacchetti aumentando il TTL di uno ogni volta: ogni pacchetto scade al router
          successivo lungo il percorso, che invia il proprio messaggio ICMP.
    \item Raccogliendo gli indirizzi sorgente dei messaggi ICMP, traceroute identifica tutti i router tra l'host locale e la
          destinazione, creando una \textbf{mappa del percorso}.
\end{itemize}

\begin{figure}[H]
\centering
\begin{tikzpicture}[every node/.style={font=\scriptsize\sffamily}]
  \node (h) at (0,0) {\icon[1cm]{laptop}}; \node[below=-2pt of h] {host};
  \node (r1) at (3.2,0) {\icon[0.9cm]{router}}; \node[below=0pt of r1] {R1};
  \node (r2) at (6.4,0) {\icon[0.9cm]{router}}; \node[below=0pt of r2] {R2};
  \node (r3) at (9.6,0) {\icon[0.9cm]{router}}; \node[below=0pt of r3] {R3};
  \node (d) at (12.8,0) {\icon[0.6cm]{server}}; \node[below=-2pt of d] {destinazione};
  \foreach \a/\b in {h/r1, r1/r2, r2/r3, r3/d} \draw[wired] (\a) -- (\b);
  % TTL 1
  \draw[msg, netblue] (0.6,0.8) -- node[above] {UDP, TTL=1} (3.0,0.8);
  \draw[msg, netred, dashed] (3.0,1.1) -- node[above] {ICMP time exceeded da R1} (0.6,1.1);
  % TTL 2
  \draw[msg, netblue] (0.6,1.7) -- node[above, pos=0.3] {UDP, TTL=2} (6.2,1.7);
  \draw[msg, netred, dashed] (6.2,2.0) -- node[above, pos=0.6] {ICMP time exceeded da R2} (0.6,2.0);
  % TTL 4
  \draw[msg, netblue] (0.6,-1.2) -- node[below, pos=0.2] {UDP, TTL=4, porta alta casuale} (12.6,-1.2);
  \draw[msg, netgreen!60!black, dashed] (12.6,-1.6) -- node[below, pos=0.6] {ICMP port unreachable: arrivato, traceroute si ferma} (0.6,-1.6);
\end{tikzpicture}
\caption{traceroute aumenta il TTL a ogni tentativo: ogni router che scarta il pacchetto si rivela con un messaggio ICMP.}
\end{figure}

\subsection{Come termina traceroute}

\begin{itemize}
    \item Il processo continua finché si raggiunge la destinazione o il \textbf{numero massimo di hop} configurato (di default 30).
    \item I pacchetti inviati da traceroute (su Linux, datagrammi UDP) hanno come porta di destinazione un \textbf{numero alto scelto a caso}.
          Se uno di questi pacchetti arriva al destinatario, di norma sul destinatario non c'è alcun processo su quella porta,
          e il destinatario risponde con un messaggio ICMP \textbf{Destination Unreachable (Port Unreachable)}.
    \item Quando il client riceve questo messaggio, traceroute termina e non invia altri segmenti UDP.
\end{itemize}

\info{Un'applicazione che scavalca il trasporto}{
    Anche se traceroute è ovviamente un'applicazione, non usa il livello di trasporto in modo diretto: prepara e invia
    direttamente pacchetti a livello di rete (IP, ICMP). Su Windows il comando equivalente, \texttt{tracert}, usa echo request
    ICMP invece di UDP.
}

\subsection{L'output e gli asterischi}

L'output è una lista ordinata di \emph{hop}, i router attraversati, con indirizzo IP o nome e i tempi di risposta
misurati per ciascun salto (di solito tre tentativi per hop).

\begin{lstlisting}[language=bash]
$ traceroute italia.it
traceroute to italia.it (151.101.195.10), 30 hops max, 60 byte packets
 1  _gateway (192.168.1.1)  1.210 ms  1.107 ms  1.049 ms
 2  100.75.5.21 (100.75.5.21)  8.921 ms  8.873 ms  8.826 ms
 3  ...
 6  151.101.195.10  25.4 ms  25.1 ms  24.9 ms
 7  * * *
 8  * * *
\end{lstlisting}

\warning{Gli asterischi}{
    Se un router non risponde in tempo, al suo posto compare un asterisco (\texttt{*}) per quel tentativo. Può indicare
    congestione, filtri, o dispositivi configurati per non rispondere ai pacchetti ICMP. Talvolta i pacchetti non ricevono
    risposta né dal destinatario né dai router intermedi: nell'esempio, dopo 6 passi si raggiunge il nodo di
    \texttt{italia.it}, che però non risponde (probabilmente per via di filtri). traceroute continua allora a inviare
    datagrammi UDP che vengono ignorati, e che non vengono nemmeno scartati per TTL scaduto: il TTL, ricordiamo, non misura
    il tempo ma il \textbf{numero di hop}.
}

% ══════════════════════════════════════════════════════════════
\section{traceroute e la frammentazione IP}
% ══════════════════════════════════════════════════════════════

\subsection{Specificare la dimensione del datagramma}

Con il traceroute di Linux si può indicare, come secondo argomento, la \textbf{lunghezza del pacchetto} sonda inviato:

\begin{lstlisting}[language=bash]
$ traceroute italia.it          # pacchetti piccoli (dimensione di default)
$ traceroute italia.it 3000     # pacchetti da 3000 byte: verranno frammentati
$ traceroute 8.8.8.8
$ traceroute 8.8.8.8 3000
\end{lstlisting}

La dimensione comprende la parte dati che viaggia nella rete più gli header UDP e IP che la incapsulano (20 byte di
header IP, 8 di header UDP, il resto è payload). Aumentandola si inviano pacchetti più grandi, che richiedono
\textbf{frammentazione} se superano l'MTU del percorso. La dimensione si riferisce ai pacchetti sonda, non ai messaggi
ICMP ricevuti, che hanno dimensioni diverse perché contengono solo una parte del pacchetto originale che ha generato
l'errore. Serve a verificare come le reti gestiscono pacchetti di dimensioni diverse.

\subsection{Osservare la frammentazione}

\begin{esercizio}{Elementi di IPv4 con traceroute e Wireshark}
\textbf{1. Elementi di IPv4.} Si seleziona il primo segmento UDP inviato da traceroute (senza specificare la dimensione) e si
espande la sezione \emph{Internet Protocol}.
\begin{enumerate}
    \item Qual è l'indirizzo IP del proprio host? \emph{(campo Source)}
    \item Qual è il valore del TTL? \emph{(1, per il primo pacchetto)}
    \item Qual è il protocollo superiore? \emph{(17, UDP)}
    \item Quanti byte ci sono nell'header IP? \emph{(20, senza opzioni)}
    \item Quanti byte ci sono nel payload? \emph{(lunghezza totale meno 20)}
    \item Il datagramma è stato frammentato? \emph{(no: flag More Fragments a 0 e offset 0)}
\end{enumerate}
Filtro per i pacchetti inviati al server: \texttt{ip.src==<mio\_IP> \&\& ip.dst==<IP\_server> \&\& udp}.
\begin{enumerate}\setcounter{enumi}{6}
    \item Quali campi restano costanti tra i pacchetti? \emph{(versione, lunghezza dell'header, indirizzi, protocollo: la sonda è sempre la stessa)}
    \item Che andamento ha il campo \emph{Identification}? \emph{(cresce di uno a ogni datagramma)}
\end{enumerate}
Filtro per i messaggi ICMP dei router: \texttt{ip.dst==<mio\_IP> \&\& icmp}. Si osservino il protocollo superiore (1, ICMP),
gli identificatori e i TTL dei datagrammi inviati dai router.

\textbf{2. Frammentazione IP.} Si ripete con dimensione 3000 byte e filtro \texttt{ip.src==<mio\_IP> \&\& ip.dst==<IP\_server> \&\& ip}.
\begin{itemize}
    \item Il pacchetto da 3000 byte viene diviso in \textbf{tre frammenti}, con offset 0, 1480 e 2960 byte (nel campo dell'header 0, 185 e 370).
    \item Indicano la frammentazione il flag \textbf{More Fragments} e il campo \textbf{Fragment Offset}.
    \item Il primo frammento ha offset 0 e More Fragments = 1; quelli successivi hanno offset maggiore di 0; l'ultimo ha More Fragments = 0.
    \item Tutti i frammenti hanno lo \textbf{stesso identificatore} (per esempio 28264) e lo stesso TTL: una volta ricomposti a
          destinazione formeranno il payload ``artificiale'' preparato da traceroute.
\end{itemize}
\end{esercizio}

% ══════════════════════════════════════════════════════════════
\section{DNS, VPN e spoofing: un caso di studio}
% ══════════════════════════════════════════════════════════════

Durante un'esercitazione si è provato a interrogare direttamente un root server, collegati attraverso una
\textbf{VPN} (una rete privata virtuale che cifra tutto il traffico verso un server remoto).

\begin{lstlisting}[language=bash]
$ nslookup iana.org c.root-servers.net
Server:   c.root-servers.net
Non-authoritative answer:
Name:     iana.org
Address:  192.0.43.8
\end{lstlisting}

Qualcosa non torna. Un root server \textbf{non} conosce i record A dei singoli domini: dovrebbe rispondere solo con i name
server del dominio \texttt{org} (è quello che succede interrogando un altro root server senza VPN, dove ``non ci sono risposte
con record A o AAAA, ma la risposta contiene informazioni sui name server per il dominio org''). Qui invece sembra che il
root server abbia fornito direttamente l'indirizzo, con una risposta non autorevole.

\begin{figure}[H]
\centering
\begin{tikzpicture}[every node/.style={font=\scriptsize\sffamily}]
  \node (h) at (0,0) {\icon[1cm]{laptop}}; \node[below=-2pt of h] {host};
  \node (v) at (5,0) {\icon[0.7cm]{server}}; \node[below=-2pt of v, align=center] {server VPN};
  \node (dns) at (10,1.3) {\icon[0.6cm]{server}}; \node[right=0pt of dns, align=left] {DNS del provider};
  \node (root) at (10,-1.3) {\icon[0.6cm]{server}}; \node[right=0pt of root] {c.root-servers.net};
  \draw[very thick, netorange] (h) -- node[above] {tunnel cifrato} node[below] {Wireshark vede solo byte cifrati} (v);
  \draw[msg, netblue] (v) -- node[above, sloped] {risolve il nome} (dns);
  \draw[msg, dashed, black!40] (v) -- node[below, sloped] {la query non arriva mai qui} (root);
  \draw[msg, netred, very thick] (4.6,0.6) to[bend right=25] node[above] {risposta ``falsificata'' come se venisse dal root server} (0.5,0.6);
\end{tikzpicture}
\caption{Il server VPN intercetta la query DNS, la risolve da sé e risponde fingendo di essere il server interrogato.}
\end{figure}

\begin{itemize}
    \item Ovviamente Wireshark non vede il traffico da e verso la VPN, che è cifrato.
    \item Evidentemente il server VPN a cui ci si collega non solo inoltra le query DNS al name server del provider, ma si prende
          carico della risoluzione del nome e poi fa \textbf{spoofing} dei datagrammi UDP/IP, confezionando la risposta come se
          provenisse dal name server interrogato.
    \item Con un altro provider VPN lo spoofing non avviene, e la risposta è quella attesa da un root server (i name server di \texttt{org}).
\end{itemize}

\info{La lezione}{
    Il campo ``indirizzo sorgente'' di un datagramma IP \textbf{non è una prova} di chi l'abbia davvero inviato: chi controlla un
    punto del percorso può falsificarlo. È esattamente il problema dell'\textbf{IP spoofing} che affronteremo nella parte sulla
    sicurezza, e il motivo per cui servono meccanismi di autenticazione.
}

% ══════════════════════════════════════════════════════════════
\section{Riepilogo}
% ══════════════════════════════════════════════════════════════

\riepilogo{
\begin{itemize}
    \item \textbf{ICMP} trasporta messaggi di controllo e di errore del livello di rete dentro datagrammi IP (tipo, codice, checksum).
    \item \textbf{ping}: echo request (8) ed echo reply (0), misura RTT e perdite; non usa il trasporto.
    \item \textbf{traceroute}: pacchetti con TTL = 1, 2, 3, \dots; ogni router che li scarta risponde con ICMP \emph{time exceeded} (11);
          la destinazione risponde con \emph{port unreachable} (3/3) e il programma si ferma. Gli asterischi indicano assenza di risposta.
    \item Con pacchetti più grandi dell'MTU si osserva la \textbf{frammentazione}: stesso identificatore, offset crescenti, flag More Fragments.
    \item L'indirizzo sorgente IP si può falsificare (\textbf{spoofing}), come mostra il caso della VPN.
\end{itemize}
}
